\documentclass[11pt,a4paper]{article}

% ----------------------------- Packages -----------------------------
\usepackage[utf8]{inputenc}
\usepackage[T1]{fontenc}
\usepackage[margin=1in]{geometry}
\usepackage{graphicx}
\usepackage{hyperref}
\usepackage{xcolor}
\usepackage{booktabs}
\usepackage{listings}
\usepackage{enumitem}
\usepackage{titlesec}
\usepackage{fancyhdr}
\usepackage{abstract}

% ----------------------------- Styling ------------------------------
\definecolor{gold}{RGB}{184,144,47}
\definecolor{ink}{RGB}{20,24,38}
\definecolor{codebg}{RGB}{245,245,248}

\hypersetup{
  colorlinks=true,
  linkcolor=ink,
  urlcolor=gold,
  citecolor=ink
}

\titleformat{\section}{\Large\bfseries\color{ink}}{\thesection}{1em}{}
\titleformat{\subsection}{\large\bfseries\color{ink}}{\thesubsection}{1em}{}

\pagestyle{fancy}
\fancyhf{}
\rhead{AiKhbar}
\lhead{Urdu AI News Intelligence Platform}
\cfoot{\thepage}

\lstset{
  basicstyle=\ttfamily\small,
  backgroundcolor=\color{codebg},
  frame=single,
  rulecolor=\color{gray!40},
  breaklines=true,
  showstringspaces=false,
  columns=fullflexible
}

% ----------------------------- Title --------------------------------
\title{\textbf{AiKhbar}\\[0.3em]
\large An Urdu AI-Powered News Intelligence and Audio Briefing Platform}
\author{Project Report}
\date{\today}

% ====================================================================
\begin{document}

\maketitle
\thispagestyle{fancy}

\begin{abstract}
\noindent
AiKhbar is a full-stack artificial-intelligence system for the Urdu-language
news domain. Rather than a conventional news reader, it is designed as an
\emph{AI media intelligence platform}: it automatically aggregates news from
multiple Pakistani sources, classifies and clusters stories, produces unified
multi-document summaries, synthesises natural Urdu audio briefings, aggregates
public and political opinion through Retrieval-Augmented Generation (RAG),
reconstructs the historical context of major events, and supports a
conversational interface for exploring the news. This report describes the
system's objectives, architecture, AI pipelines, data model, implementation
and evaluation, and outlines a roadmap toward production scale.
\end{abstract}

\tableofcontents
\newpage

% ====================================================================
\section{Introduction}

\subsection{Motivation}
Consumers of Urdu-language news face three persistent problems: information
\emph{overload} from many overlapping sources, a lack of \emph{synthesis}
across outlets reporting the same event, and limited \emph{accessibility} for
users who prefer listening over reading. Existing aggregators present raw
headlines but do little to summarise, contextualise, or analyse coverage.

\subsection{Objectives}
AiKhbar addresses these problems with the following goals:
\begin{itemize}[noitemsep]
  \item Aggregate real-time Urdu news automatically from multiple sources.
  \item Classify articles into seven categories and cluster articles that
        describe the same story.
  \item Generate concise, faithful Urdu summaries at multiple depths.
  \item Convert summaries into natural Urdu audio briefings.
  \item Aggregate diverse opinions and reconstruct historical timelines.
  \item Provide a conversational, RAG-grounded news assistant.
  \item Deliver a personalised one-click daily brief.
\end{itemize}

\subsection{Scope}
The work delivers a complete Minimum Viable Product (MVP): an asynchronous
FastAPI backend, a React frontend, AI pipelines built on a hosted large
language model (LLM), a vector database for retrieval, and full deployment
tooling. Production hardening (authentication, distributed task processing) is
identified as future work.

% ====================================================================
\section{Literature and Background}

\subsection{Large Language Models for News}
Modern instruction-tuned LLMs perform abstractive summarisation, zero-shot
classification and question answering without task-specific training. AiKhbar
uses NVIDIA NIM, a hosted, OpenAI-API-compatible inference service, allowing
open and commercial models to be accessed through a single client.

\subsection{Retrieval-Augmented Generation}
RAG combines a parametric LLM with a non-parametric retrieval store. Source
documents are embedded into a vector space; at query time the most relevant
passages are retrieved and supplied to the LLM as grounding context. This
reduces hallucination and lets answers cite real sources --- essential for a
news product.

\subsection{Multilingual Embeddings}
AiKhbar indexes Urdu articles but may retrieve English editorials. The
\texttt{multilingual-e5-base} sentence-embedding model places both languages in
a shared semantic space, enabling cross-lingual retrieval.

% ====================================================================
\section{System Architecture}

\subsection{Design Principles}
The system follows \textbf{clean architecture}: dependencies point inward, from
the HTTP layer toward domain logic and infrastructure. Concretely:
\begin{itemize}[noitemsep]
  \item API routes are thin --- they validate input and delegate to services.
  \item Each AI capability is an isolated, independently testable package.
  \item All LLM calls pass through a single gateway module.
  \item Prompt templates are centralised, separate from business logic.
\end{itemize}

\subsection{Layered Overview}
\begin{table}[h]
\centering
\begin{tabular}{@{}lll@{}}
\toprule
\textbf{Layer} & \textbf{Package} & \textbf{Responsibility} \\
\midrule
HTTP          & \texttt{api/v1}, \texttt{websocket} & Routing, validation \\
Orchestration & \texttt{services/} & Pipeline, daily brief, scheduler \\
Domain        & \texttt{scraping}, \texttt{rag}, \texttt{tts}, etc. & AI pipelines \\
Infrastructure& \texttt{db}, \texttt{models}, \texttt{core} & Persistence, config \\
\bottomrule
\end{tabular}
\caption{The four architectural layers of AiKhbar.}
\end{table}

\subsection{Technology Stack}
The backend uses Python~3.12 with FastAPI and asynchronous SQLAlchemy~2 over
PostgreSQL. Background scheduling uses APScheduler. The frontend is a React~18
single-page application built with Vite, styled with TailwindCSS, animated with
Framer Motion, and managed with Redux Toolkit. AI components use NVIDIA NIM
(LLM), SentenceTransformers (embeddings) and FAISS (vector index).

% ====================================================================
\section{AI Pipelines}

\subsection{News Ingestion}
The ingestion pipeline executes six stages with isolated error handling:
\emph{scrape} (RSS parsing and full-text extraction), \emph{deduplicate}
(exact matching on URL and a normalised content hash), \emph{classify},
\emph{persist}, \emph{index} (into FAISS), and \emph{cluster}. A per-stage
report records counts and errors so partial failures degrade gracefully.

\subsection{Classification}
Articles are classified into Politics, Sports, Economy, Crime, International,
Technology and Entertainment. The MVP uses the LLM as a zero-shot classifier:
it handles Urdu without training data and returns category, confidence, tags
and region in a single structured call. The interface allows a fine-tuned
transformer classifier to replace it later.

\subsection{Story Clustering and Summarisation}
Articles from different outlets covering the same event are grouped using
agglomerative clustering with a cosine-distance threshold on headline-and-lead
embeddings, performed per category to avoid cross-topic merges. Each cluster is
summarised by the LLM into one \emph{unified} Urdu summary, available at four
depths: short, detailed, a five-minute morning brief, and a fifteen-minute deep
dive.

\subsection{Retrieval-Augmented Generation}
Article bodies are split into overlapping 220-word chunks, embedded, and stored
in a FAISS \texttt{IndexFlatIP} index (inner product over normalised vectors
yields cosine similarity). A single retriever serves three features:
conversational chat, opinion aggregation and timeline reconstruction. Retrieved
context is formatted and injected into the LLM prompt.

\subsection{Opinion Aggregation and Timeline}
For a story, the opinion aggregator retrieves a broad set of related passages
and prompts the LLM to identify distinct perspectives --- public sentiment,
pro-government, opposition, expert and international --- each with a stance and
confidence. The timeline builder similarly reconstructs a chronological
sequence of background events.

\subsection{Media Analysis}
The media-analysis engine assesses each article's tone, political leaning,
sentiment polarity and propaganda indicators, persisting the result so the
interface can present a media-bias badge.

\subsection{Urdu Audio Generation}
Summaries are converted into broadcast-style narration scripts with an
introduction, headline emphasis and inter-story pauses. A provider abstraction
selects the Urdu Orator API as primary and Coqui XTTS-v2 as an open-source
fallback; synthesised audio is cached on disk keyed by content hash.

\subsection{Personalisation}
A deliberately simple and explainable model learns per-category affinity as a
time-decayed, action-weighted sum of user interactions. The personalised feed
ranks recent articles by a weighted combination of affinity and recency.

% ====================================================================
\section{Data Model}

The relational schema centres on the \texttt{story\_clusters} entity. An
\texttt{Article} belongs to at most one cluster; a cluster owns its
\texttt{Summary}, \texttt{Opinion} and \texttt{TimelineEvent} records. A
\texttt{User} has many \texttt{UserInteraction} records that drive
personalisation. All tables carry creation and update timestamps. The complete
schema is documented in \texttt{docs/DATABASE.md}.

% ====================================================================
\section{Implementation}

\subsection{Backend}
The backend exposes a versioned REST API plus WebSocket endpoints for live
chat. CPU-bound work --- embedding, FAISS operations and local TTS inference
--- is offloaded from the event loop with \texttt{asyncio.to\_thread}. Periodic
ingestion is driven by APScheduler with overlap protection.

\subsection{Frontend}
The frontend implements a cinematic, editorial ``dark luxury'' aesthetic using
glassmorphism, layered animated backgrounds and reusable motion components.
Pages include a landing page, a filterable newsroom, a story detail view with
on-demand opinion/timeline/audio enrichment, the daily brief, and a
conversational chat interface.

\subsection{Reliability}
The system adopts a \emph{real-calls-only} policy: missing configuration (for
example, an absent API key) raises an explicit error rather than silently
degrading. External calls use exponential-backoff retries, and each pipeline
stage isolates its failures.

% ====================================================================
\section{Evaluation}

\subsection{Functional Verification}
All backend modules compile and the core configuration, ORM models, schemas and
scraping registry import successfully. The unit-test suite (text chunking,
content hashing, narration assembly) passes. The frontend builds cleanly into
an optimised production bundle.

\subsection{Qualitative Assessment}
The architecture achieves strong separation of concerns: the LLM provider, the
vector index implementation and the TTS provider can each be replaced by
editing a single module. The pipeline's per-stage error isolation makes the
system resilient to upstream degradation.

\subsection{Limitations}
The MVP runs background work in-process rather than via a distributed queue;
the FAISS flat index is exact but does not scale to very large corpora; and
classification depends on LLM availability. Each limitation has a planned
remedy in the roadmap.

% ====================================================================
\section{Future Work}

Planned enhancements include: JWT authentication and rate limiting; migration
to Celery and Redis for distributed task processing; an IVF/HNSW or managed
vector index; token-streaming chat responses; a fine-tuned Urdu classifier;
cross-lingual ingestion; and a mobile client. The full phased plan is given in
\texttt{docs/ROADMAP.md}.

% ====================================================================
\section{Conclusion}

AiKhbar demonstrates that a modular, clean-architecture design can integrate
multiple AI capabilities --- aggregation, classification, summarisation,
retrieval-augmented analysis, speech synthesis and personalisation --- into a
coherent media-intelligence platform for an under-served language. The system
is functional end to end, maintainable, and structured so that each component
can evolve independently toward production scale.

% ====================================================================
\appendix
\section{Key API Endpoints}
\begin{table}[h]
\centering
\begin{tabular}{@{}lll@{}}
\toprule
\textbf{Method} & \textbf{Path} & \textbf{Purpose} \\
\midrule
GET  & \texttt{/api/v1/news} & List / filter articles \\
POST & \texttt{/api/v1/news/scrape} & Trigger ingestion \\
GET  & \texttt{/api/v1/stories/top} & Top story clusters \\
POST & \texttt{/api/v1/stories/\{id\}/summary} & Unified summary \\
POST & \texttt{/api/v1/stories/\{id\}/opinions} & RAG opinion aggregation \\
POST & \texttt{/api/v1/briefs/daily} & One-click daily brief \\
POST & \texttt{/api/v1/chat} & Conversational Q\&A \\
POST & \texttt{/api/v1/tts/synthesize} & Urdu text-to-speech \\
\bottomrule
\end{tabular}
\caption{Selected REST endpoints.}
\end{table}

\section{Repository}
Source code: \url{https://github.com/omerrfarooqq/AiKhbar}

\end{document}
